% Part: incompleteness
% Chapter: introduction
% Section: overview

\documentclass[../../../include/open-logic-section]{subfiles}

\begin{document}

\olfileid{inc}{int}{ovr}

\olsection{अपूर्णतेच्या निष्कर्षांचा आढावा}

गणिताला स्वयंसिद्धकांनी निरूपणयोग्य !!{axiomatizable} उपपत्तीत
आकारिक रूप देता येईल आणि ती !!{complete} व !!{decidable} आहे हे
सिद्ध करता येईल, अशी हिल्बर्टची अपेक्षा होती. शिवाय, अतिशय दुर्बल,
``सांतवादी'' साधनांनी त्या उपपत्तीची सुसंगतता सिद्ध करून
अंतःप्रज्ञावादाच्या आक्षेपांपासून अभिजात गणिताचे समर्थन करण्याचे
त्याचे ध्येय होते. गोडेलच्या अपूर्णता प्रमेयांनी ही उद्दिष्टे
साध्य होऊ शकत नाहीत हे दाखवले.

गोडेलच्या पहिल्या अपूर्णता प्रमेयाने रसेल आणि व्हाइटहेड यांच्या
\emph{Principia Mathematica} ची एक आवृत्ती !!{complete} नाही हे
दाखवले. पण सिद्धता प्रत्यक्षात अतिशय व्यापक होती आणि अनेक प्रकारच्या
उपपत्तींना लागू होते. म्हणजे केवळ \emph{Principia Mathematica}
गणित पूर्णपणे पकडू शकला नाही असे नव्हे, तर कोणतीही \emph{स्वीकारार्ह}
उपपत्ती तसे करू शकत नाही. अपूर्णता प्रमेये लागू होण्यासाठी उपपत्तीत
कोणते गुणधर्म पुरेसे आहेत हे वेगळे ओळखण्यास आणि गोडेलची सिद्धता
केवळ त्यांवर अवलंबून राहील अशी व्यापक करण्यास काही काळ लागला.
आता मात्र अंकगणिताच्या $\Lang L_A$ भाषेतील उपपत्त्यांसाठी पहिले
अपूर्णता प्रमेय अतिशय व्यापक रूपात मांडता येते.

\begin{thm}
$\Gamma$ ही~$\Lang
L_A$ मधील सुसंगत आणि !!{axiomatizable} उपपत्ती असेल, आणि ती सर्व
संगणनक्षम फलनांचे व !!{decidable} संबंधांचे !!{represents} करत असेल,
तर $\Gamma$ !!{complete} नाही.
\end{thm}

$\Gamma$ !!{complete} नाही याचा अर्थ किमान एक असे
!!{sentence}~$!A$ आहे की $\Gamma \Proves/ !A$ आणि
$\Gamma \Proves/ \lnot
!A$. अशा !!a{sentence} ला ($\Gamma)$ पासून \emph{स्वतंत्र}
म्हणतात. अशी स्वतंत्र वाक्ये असलीच पाहिजेत हे प्रत्यक्षात तुलनेने
लवकर सिद्ध करता येते. परंतु गोडेलच्या सिद्धतेचे सामर्थ्य असे की
ती स्वतःच्या असिद्धतायोग्यतेचा दावा करणाऱ्या !!{sentence} चे
\emph{विशिष्ट उदाहरण} देते.
या रंजक रचनेत~$\Gamma$ साठीचे \emph{गोडेल-वाक्य} म्हणून ओळखले
जाणारे !!a{sentence}~$!G_\Gamma$ तयार होते. ते सिद्ध करता येत
नाही, कारण $\Gamma$ मध्ये $!G_\Gamma$ हेच वाक्य~$\Gamma$ मध्ये
$!G_\Gamma$ सिद्ध करता येत नाही या दाव्याशी सममूल्य असते.
ही रचना \emph{रचनात्मकपणे} केली जाते: $\Gamma$ ची स्वयंसिद्धकीय
मांडणी आणि !!{derivation} पद्धतीचे वर्णन दिल्यास सिद्धतेतून
$!G_\Gamma$ प्रत्यक्ष लिहिण्याची पद्धत मिळते. केवळ सुसंगततेवरून हे मूळ गोडेल-वाक्य
असिद्धतायोग्य असते; त्याचे खंडनही अशक्य ठरवण्यासाठी अधिक प्रबळ
गृहीतक, उदाहरणार्थ ओमेगा-सुसंगतता, आवश्यक असते. मात्र रॉसरच्या
बदललेल्या रचनेत केवळ सुसंगततेवरून स्वतंत्र वाक्य मिळते; वरील
प्रमेय त्या अधिक व्यापक रूपात दिले आहे.
%READERNOTE{SOL6-A056 — गोठवलेल्या इंग्रजी आढाव्यात मूळ Gödel sentence आणि केवळ consistency गृहीत धरणाऱ्या व्यापक प्रमेयाची गल्लत होते. मूळ वाक्याच्या स्वतंत्रतेसाठी अतिरिक्त गृहीतक आणि Rosser च्या बदललेल्या रचनेचा फरक येथे स्पष्ट केला.}

गोडेलच्या सिद्धतेतील रचनेसाठी $\Lang L_A$ मधील पदे व
!!{formula}s यांचे गुणधर्म आणि त्यांवरील क्रिया त्याच $\Lang L_A$
मध्ये व्यक्त करण्याचा मार्ग हवा. यात ``$!A$ हे !!a{sentence} आहे'',
``$\delta$ ही~$!A$ ची !!a{derivation} आहे'' यांसारखे गुणधर्म
आणि $\Subst{!A}{t}{x}$ सारख्या क्रिया येतात. या मार्गाने
(a) चिन्हे आणि त्यांचे अनुक्रम (म्हणजे पदे, !!{formula}s,
!!{derivation}s इत्यादी) नैसर्गिक संख्यांत सांकेतिक करून
त्या गुणधर्मांचे व संबंधांचे निरूपण करावे लागते, कारण $\Lang L_A$
नैसर्गिक संख्यांविषयी बोलू शकते. तसेच (b) $\Gamma$ ला संबंधित
तथ्ये सिद्ध करता येतील अशा रीतीने हे करावे लागते. म्हणून हे
गुणधर्म नैसर्गिक संख्यांच्या !!{decidable} गुणधर्मांनी आणि क्रिया
त्यांवरील संगणनक्षम फलनांनी सांकेतिक होतात हे दाखवावे लागते.
याला ``विन्यासमीमांसेचे अंकगणितीकरण'' म्हणतात.

विन्यासमीमांसेचे अंकगणितीकरण कसे करावे हे पाहण्यापूर्वी, $\Gamma$
``पुरेशी प्रबळ'' आहे या अटीचा विचार करू; म्हणजे ती सर्व संगणनक्षम
फलनांचे आणि !!{decidable} संबंधांचे !!{represents} करते. यासाठी
``संगणनक्षम'' याची नेमकी व्याख्या द्यावी लागेल. ती अनेक मार्गांनी
देता येते; उदाहरणार्थ, ट्यूरिंग यंत्रांच्या प्रतिमानातून किंवा
सर्वसाधारण वापराच्या संगणक-भाषेतील कार्यक्रमांनी संगणित होणाऱ्या
फलनांच्या रूपात. मात्र आपले उद्दिष्ट या फलनांचे आणि संबंधांचे
भाषा~$\Lang L_A$ मधील उपपत्तीत !!{represents}s करणे असल्याने
केवळ संख्यात्मक फलनांच्या संगणनक्षमतेची सोपी व्याख्या निवडणे
उत्तम. ही \emph{पुनरावर्ती फलन} ही संकल्पना आहे. म्हणून प्रथम
पुनरावर्ती फलनांची चर्चा करू. मग $\Th{Q}$ आधीच सर्व पुनरावर्ती
फलनांचे आणि संबंधांचे !!{represents} करते हे दाखवू. त्यामुळे $\Th{Q}$
आणि $\Th{PA}$ यांसारख्या विशिष्ट उपपत्त्यांना अपूर्णता प्रमेय
लागू करता येईल, कारण त्या ``पुरेशा प्रबळ'' उपपत्ती असल्याचे
आपण प्रस्थापित केलेले असेल.

विन्यासमीमांसेच्या अंकगणितीकरणातून अखेरीस !!a{formula}
$\OProv[\Gamma](x)$ मिळते. !!{formula}s ना संख्यांनी सांकेतिक
करण्यामुळे हे सूत्र~$\Gamma$ च्या स्वयंसिद्धकांपासून सिद्धतायोग्यता
व्यक्त करते. विशेषतः, $!A$ चा संकेतांक~$n$ असेल आणि $\Gamma \Proves
!A$ असेल, तर $\Gamma \Proves \Prov[\Gamma](\num{n})$.
$\Gamma$ साठीचे हे ``सिद्धतायोग्यता विधेय'' विशिष्ट अर्थाने
$\Gamma$ ची सुसंगतताही~$\Lang L_A$ मधील !!a{sentence} म्हणून
व्यक्त करू देते. $\Gamma$ साठीचे ``सुसंगतता-विधान'' हे
$\lnot \Prov[\Gamma](\num{n})$ हे !!{sentence} मानू, जिथे $n$ हा
एखाद्या विरोधाभासाचा, उदा. ~ $\lfalse$ चा, संकेतांक आहे.
दुसरे अपूर्णता प्रमेय सांगते की सुसंगत !!{axiomatizable} उपपत्ती
स्वतःची सुसंगतता-विधानेही सिद्ध करू शकत नाहीत. हे प्रमेय लागू
होण्यासाठी लागणाऱ्या अटी, उपपत्ती सर्व संगणनक्षम फलने आणि
!!{decidable} संबंध निरूपित करते एवढ्यापेक्षा काहीशा अधिक
कडक आहेत; पण $\Th{PA}$ त्या पूर्ण करते हे आपण दाखवू.

\end{document}
